\documentclass[11pt]{article}
\usepackage[margin=1in]{geometry}
\usepackage{amsmath,amssymb,amsthm,longtable,booktabs,array}
\usepackage[T1]{fontenc}
\usepackage{lmodern}
\usepackage[colorlinks=true,urlcolor=blue,linkcolor=blue,citecolor=blue]{hyperref}
\newtheorem{theorem}{Theorem}
\newtheorem{lemma}{Lemma}
\newcommand{\C}{\mathcal C}
\newcommand{\B}{\mathcal B}
\newcommand{\R}{\mathbb R}
\title{A finite certificate against the choice-map characterization\newline of nonideal clutters}
\author{Draft research manuscript; authorship not assigned}
\date{31 October 2026 (America/Los\_Angeles)}
\begin{document}
\maketitle
\begin{center}\textbf{DRAFT / CANDIDATE 388087. Not submitted or accepted.}\end{center}
\begin{abstract}
We describe an explicit clutter of eleven triples on eleven elements whose real covering polyhedron has a fractional extreme point, but which admits no pair of choice maps satisfying Kir\'aly's proposed intersection condition. Nonideality follows from a short linear elimination. Nonexistence of the maps is proved by a finite branching certificate with 6,517 leaves, verified by a checker whose soundness is proved in Lean. The exact theorem file, data, pinned dependency manifest and replay instructions accompany this draft. The underlying small cubic Lehman object belongs to an already studied class; neither its construction nor any priority or novelty claim is asserted here. Independent fidelity review has passed, and a separate significance review gives a qualified bounded-prior PASS. This draft is not accepted and makes no worldwide-priority claim.
\end{abstract}

\section{Question, sources and precise scope}
A clutter on a finite set $V$ is a family $\C$ of subsets with no proper containment between members. Its blocker $b(\C)$ is the family of inclusion-minimal transversals: sets meeting every member of $\C$. Set
\[
 P(\C)=\{x\in\R^V:x_v\geq0\ (v\in V),\quad
                         \sum_{v\in X}x_v\geq1\ (X\in\C)\}.
\]
A clutter is ideal if its covering polyhedron is integral. Here we use the vertex formulation: every extreme point has integer coordinates. Our nonideality proof gives an explicit noninteger extreme point of this actual real polyhedron, so does not depend on a surrogate finite test for ideality. In particular $P(\C)$ is pointed since it is contained in the nonnegative orthant.

Define the choice-map property $K(\C)$ by the existence of
\[
 p:\C\longrightarrow V,\qquad q:b(\C)\longrightarrow V
\]
with, for every $X\in\C$, $Y\in b(\C)$,
\begin{equation}\label{eq:choices}
 p(X)\in X,\quad q(Y)\in Y,\quad
 p(X)=q(Y)\ \Longrightarrow\ |X\cap Y|>1.
\end{equation}
The intended MathDB 388087/Egres question attributed to Tam\'as Kir\'aly is
\begin{equation}\label{eq:conj}
 \text{for every finite clutter }\C,\qquad
 \C\text{ is nonideal}\quad\Longleftrightarrow\quad K(\C).
\end{equation}
Pap's January 2011 revision \cite[Conjecture 1.2, p.~1]{PapRevised} directly states the separate-$p,q$ form, with both families nonideal on the left. Pap's November 2010 report \cite[Definition 1.1 and Conjecture 1.2, p.~1]{PapTR} formulates the nonideal-to-choice direction using a single function $h:\C\cup b(\C)\to V$, chosen inside each set. Restricting $h$ to the two families gives $p,q$ satisfying \eqref{eq:choices}. Consequently nonexistence of $p,q$ also excludes this single-$h$ form; equivalence of the two forms for arbitrary overlapping families is not needed.

Pap's thesis \cite[Conjecture 3.4.5, pp.~70--71]{PapThesis}, also checked in the author-hosted copy listed in the bibliography, explicitly uses two maps and $q(Y)\in Y$. Its left side says that both the clutter and its blocker are nonideal, in accordance with the standard blocker-invariance of ideality recalled in \cite[p.~1]{PapTR}. The formal endpoint below uses the Egres/MathDB left side, nonideality of $\C$. It does not formalize the general blocker-invariance theorem. The Egres snapshot \cite{Egres} and MathDB wording have $p(Y)\in Y$ in their second clause; this is a domain typo, resolved directly by the January 2011 report and the thesis's explicit clause (ii), also consistent with the original report's single-$h$ formulation. Correcting that typo is not a mathematical contribution claimed by this draft.

\begin{theorem}\label{thm:main}
On $V=\{0,1,\ldots,10\}$, let $\C_{11}=\{X_0,\ldots,X_{10}\}$, where the triples in index order are
\[
\begin{array}{c|l@{\qquad}c|l}
0&\{1,7,10\}&1&\{2,3,8\}\\
2&\{2,3,4\}&3&\{0,3,4\}\\
4&\{0,1,4\}&5&\{0,5,6\}\\
6&\{5,6,7\}&7&\{1,6,7\}\\
8&\{2,8,9\}&9&\{8,9,10\}\\
10&\{5,9,10\}&&
\end{array}
\]
Then $\C_{11}$ is a clutter, $x^*=(1/3,\ldots,1/3)$ is a fractional extreme point of $P(\C_{11})$, and $K(\C_{11})$ is false. Thus \eqref{eq:conj} fails in the nonideal-to-choice direction.
\end{theorem}
The triples are distinct and of equal size, so the clutter assertion is immediate. The rest of the proof occupies Sections 2--4.

\paragraph{Attribution of the object and limits.}
The object is a cubic $(11,3,4)$ Lehman clutter, a member of a known small class. Mayhew, Pivotto and Royle \cite[Definition 1.1, Section 5 and Table 1, pp.~25--26]{MPR} define the type and catalogue the four classes of this type, discussing their correspondence with the earlier L\"utolf--Margot catalogue. We attribute the underlying object to this existing catalogue tradition, and do not claim a new clutter or a new classification. For a directly checkable type identification, the accompanying \texttt{logs/data-check.json} gives eleven four-sets $Z_i$ in edge order with $|X_j\cap Z_i|=1+\delta_{ij}$; Explicitly, using the Table~\ref{tab:witness} indices,
\[
(Z_0,\ldots,Z_{10})=(Y_{27},Y_1,Y_{26},Y_{13},Y_2,Y_{24},Y_5,Y_{10},Y_{12},Y_{31},Y_7).
\]
Row and column degrees of the triple matrix are three. Hence its incidence matrix $A$ has a mate with $AB^{\mathsf T}=J+I$. This identification is supplementary exact arithmetic, not needed in the Lean refutation. No catalogue index or canonical isomorphism label is asserted. No minimal-nonideality, smallest-counterexample or global-priority claim is made.

\section{A fractional vertex by linear elimination}
Each edge has three elements, so $x^*$ is feasible and all eleven covering inequalities are tight. Suppose $y$ satisfies $\sum_{v\in X_i}y_v=1$ for every $i$. Subtracting equations in the pairs
\[
 (1,2),(2,3),(3,4),(5,6),(6,7),(8,9),(9,10)
\]
gives, respectively,
\[
 y_8=y_4,\quad y_2=y_0,\quad y_3=y_1,\quad
 y_0=y_7,\quad y_5=y_1,\quad y_2=y_{10},\quad y_8=y_5.
\]
Thus write $y_0=y_2=y_7=y_{10}=a$, $y_1=y_3=y_4=y_5=y_8=b$, $y_6=c$, $y_9=d$. Equations $0,2,5,8$ now read
\[
2a+b=1,\qquad a+2b=1,\qquad a+b+c=1,\qquad a+b+d=1.
\]
They force $a=b=c=d=1/3$. This proves uniqueness of the tight solution with ordinary real arithmetic, without numerical rank estimates.

If $x^*=\lambda y+(1-\lambda)z$ for feasible $y,z$ and $0<\lambda<1$, each edge sum at $y$ and $z$ is at least one, while their strictly weighted average is one. Both sums therefore equal one. Uniqueness gives $y=z=x^*$, proving extremality. Since $0<1/3<1$ there is no integer coordinate value, and $\C_{11}$ is nonideal. In Lean these steps are \texttt{eq\_xstar}, \texttt{xstar\_extreme} and \texttt{C11\_not\_ideal}.

For completeness, the thesis's literal ``both families nonideal'' premise can also be checked for this instance without appealing to general blocker duality. Each ground element is in three triples. Any transversal $Y$ therefore satisfies $11\leq\sum_i|X_i\cap Y|=3|Y|$, so $|Y|\geq4$. Hence $w^*=(1/4,\ldots,1/4)$ is feasible for $P(b(\C_{11}))$. Let $B$ have the eleven four-sets $Z_i$ mentioned above as rows. The checked identity $AB^{\mathsf T}=I+J$ makes $B$ nonsingular: $I+J$ is nonsingular since $(I+J)u=0$ forces $u_i=-\sum_j u_j$ and then $12\sum_j u_j=0$. The eleven tight constraints from $Z_i$ uniquely determine $w^*$, which is thus a fractional extreme point by the same convex-combination argument. This additional ordinary derivation uses only blocker memberships and the explicit mate identity; it is not an additional Lean endpoint or a separate credited outcome.

\section{Minimal transversals and the obstruction}
Table~\ref{tab:witness} lists the 36 transversals used in the certificate. Each listed $Y_j$ meets all eleven triples, which can be checked directly from the displayed lists. For each $v\in Y_j$, the last column gives an index $i$ with $X_i\cap Y_j=\{v\}$. Removing $v$ therefore destroys transversality; any proper subset omits some $v$, and also misses that witnessing edge. This proves $Y_j\in b(\C_{11})$ for every listed $j$. Lean proves precisely these facts as \texttt{transversal\_Y}, \texttt{minimal\_Y} and \texttt{Y\_mem\_blocker}.

\begin{longtable}{rll}
\caption{Sets used and minimality witnesses. A pair $v:i$ means $X_i\cap Y_j=\{v\}$.}\label{tab:witness}\\
\toprule $j$&$Y_j$&Witnesses $v:i$\\\midrule\endfirsthead
\toprule $j$&$Y_j$&Witnesses $v:i$\\\midrule\endhead
0 & $\{0,1,2,5,8\}$ & \texttt{0:3, 1:0, 2:2, 5:6, 8:9} \\
1 & $\{1,3,5,8\}$ & \texttt{1:0, 3:2, 5:5, 8:8} \\
2 & $\{1,4,5,8\}$ & \texttt{1:0, 4:2, 5:5, 8:1} \\
3 & $\{0,2,5,7,8\}$ & \texttt{0:3, 2:2, 5:10, 7:0, 8:9} \\
4 & $\{0,3,5,7,8\}$ & \texttt{0:4, 3:2, 5:10, 7:0, 8:8} \\
5 & $\{4,5,7,8\}$ & \texttt{4:2, 5:5, 7:0, 8:1} \\
6 & $\{0,1,2,5,9\}$ & \texttt{0:3, 1:0, 2:1, 5:6, 9:9} \\
7 & $\{1,3,5,9\}$ & \texttt{1:0, 3:1, 5:5, 9:8} \\
8 & $\{1,2,4,5,9\}$ & \texttt{1:0, 2:1, 4:3, 5:5, 9:9} \\
9 & $\{0,1,2,6,9\}$ & \texttt{0:3, 1:0, 2:1, 6:6, 9:9} \\
10 & $\{1,3,6,9\}$ & \texttt{1:0, 3:1, 6:5, 9:8} \\
11 & $\{1,2,4,6,9\}$ & \texttt{1:0, 2:1, 4:3, 6:5, 9:9} \\
12 & $\{0,2,7,9\}$ & \texttt{0:3, 2:1, 7:0, 9:9} \\
13 & $\{0,3,7,9\}$ & \texttt{0:4, 3:1, 7:0, 9:8} \\
14 & $\{2,4,5,7,9\}$ & \texttt{2:1, 4:3, 5:5, 7:0, 9:9} \\
15 & $\{3,4,5,7,9\}$ & \texttt{3:1, 4:4, 5:5, 7:0, 9:8} \\
16 & $\{2,4,6,7,9\}$ & \texttt{2:1, 4:3, 6:5, 7:0, 9:9} \\
17 & $\{3,4,6,7,9\}$ & \texttt{3:1, 4:4, 6:5, 7:0, 9:8} \\
18 & $\{1,4,6,8,9\}$ & \texttt{1:0, 4:2, 6:5, 8:1, 9:10} \\
19 & $\{0,4,7,8,9\}$ & \texttt{0:5, 4:2, 7:0, 8:1, 9:10} \\
20 & $\{4,6,7,8,9\}$ & \texttt{4:2, 6:5, 7:0, 8:1, 9:10} \\
21 & $\{0,1,2,5,10\}$ & \texttt{0:3, 1:7, 2:1, 5:6, 10:9} \\
22 & $\{1,2,3,5,10\}$ & \texttt{1:4, 2:8, 3:3, 5:5, 10:9} \\
23 & $\{1,2,4,5,10\}$ & \texttt{1:7, 2:1, 4:3, 5:5, 10:9} \\
24 & $\{0,2,6,10\}$ & \texttt{0:3, 2:1, 6:6, 10:0} \\
25 & $\{1,2,3,6,10\}$ & \texttt{1:4, 2:8, 3:3, 6:5, 10:9} \\
26 & $\{2,4,6,10\}$ & \texttt{2:1, 4:3, 6:5, 10:0} \\
27 & $\{0,2,7,10\}$ & \texttt{0:3, 2:1, 7:6, 10:9} \\
28 & $\{2,4,5,7,10\}$ & \texttt{2:1, 4:3, 5:5, 7:7, 10:9} \\
29 & $\{0,3,6,8,10\}$ & \texttt{0:4, 3:2, 6:6, 8:8, 10:0} \\
30 & $\{1,3,6,8,10\}$ & \texttt{1:4, 3:2, 6:5, 8:8, 10:10} \\
31 & $\{4,6,8,10\}$ & \texttt{4:2, 6:5, 8:1, 10:0} \\
32 & $\{0,3,7,8,10\}$ & \texttt{0:4, 3:2, 7:6, 8:8, 10:10} \\
33 & $\{0,4,7,8,10\}$ & \texttt{0:5, 4:2, 7:6, 8:1, 10:10} \\
34 & $\{0,3,6,9,10\}$ & \texttt{0:4, 3:1, 6:6, 9:8, 10:0} \\
35 & $\{3,4,6,9,10\}$ & \texttt{3:1, 4:4, 6:5, 9:8, 10:0} \\
\bottomrule\end{longtable}

Only membership of this list in the blocker is required: a proposed map on the full blocker restricts to all listed sets. Equality of this list with the whole blocker, although confirmed by an independent finite review, is not a Lean theorem used here. Nor does the proof need minimal nonideality.

Fix any $p$ choosing one element of each triple. For each $j$ and $v\in Y_j$, define
\[
 I(j,v)=\{i\in\{0,\ldots,10\}:X_i\cap Y_j=\{v\}\}.
\]
Call $v$ forbidden for $Y_j$ if $p(X_i)=v$ for some $i\in I(j,v)$. If every $v\in Y_j$ is forbidden, no value of $q(Y_j)$ is possible: its value must belong to $Y_j$, and the witnessing edge violates \eqref{eq:choices}. The certificate establishes that some listed set has every element forbidden for every possible $p$.

\section{The finite certificate is part of the proof}
The exact certificate is the literal value \texttt{cert} in the accompanying Lean file; a human-readable JSON encoding is \texttt{certificate.json}. It is a ternary tree with 3,258 internal nodes, 6,517 leaves and maximum depth 11. Its leaves use every index $0,\ldots,35$. We do not assert that this is a short hand proof, or that merely observing an exhaustive program run proves the theorem.

Here is the checker's full mathematical specification. Let $D_i$ be the three elements of $X_i$ in the displayed order, and let $K_j$ contain the pairs $(v,I(j,v))$ for $v\in Y_j$. A partial assignment $a$ is a list of pairs $(i,v)$ saying $p(X_i)=v$. A leaf labelled $j$ passes exactly when $j<36$ and for every $(v,I)\in K_j$ there is an $i\in I$ with $(i,v)\in a$. A node labelled $i$ passes exactly when $D_i$ is a three-element list $[d_0,d_1,d_2]$ and all three child checks pass under the extensions $(i,d_k)::a$, $k=0,1,2$. Otherwise it fails.

\begin{lemma}[Soundness of the checker]
If a tree passes under $a$, every total assignment $b$ consistent with $a$ and choosing $b(i)\in D_i$ reaches a listed $j$ such that for every $(v,I)\in K_j$ some $i\in I$ satisfies $b(i)=v$.
\end{lemma}
\begin{proof}
Induct on the tree. At a leaf every required pair is already in $a$, so consistency supplies the claimed equalities. At a node, $b(i)$ equals one of the three listed values. The corresponding child passes and $b$ is consistent with its extended assignment; apply the induction hypothesis. This does not assume that different branches correspond to different globally consistent assignments.
\end{proof}
This lemma is \texttt{chk\_sound}. Lean checks the literal finite data by \texttt{decide +kernel}: \texttt{cert\_ok} proves the root check succeeds on the empty assignment. The data bridges \texttt{fact\_dom}, \texttt{fact\_cover} and \texttt{fact\_single} connect the natural-number lists to the finite-set definitions. Specifically they ensure every actual choice lies in its domain list, every possible $q(Y_j)$ has an entry in $K_j$, and every associated edge index is in range and has intersection cardinality at most one. These weaker bridge facts already suffice for the contradiction, even without using the full exactness of $I(j,v)$.

Given $p,q$, apply soundness to $b(i)=p(X_i)$, select the blocked $Y_j$, and use the entry corresponding to $q(Y_j)$. It supplies an edge $i$ with $p(X_i)=q(Y_j)$ and $|X_i\cap Y_j|\leq1$, contrary to \eqref{eq:choices}. This is \texttt{not\_kiraly}. Together with Section 2 it proves Theorem~\ref{thm:main}.

\section{Formal endpoint, reproducibility and trust scope}
All formal definitions and proofs are in the unchanged file
\texttt{KiralyCounterexample.lean}. Its SHA-256 is
\begin{quote}\small\ttfamily
1370129f56ef40e2dc317a8cf12ef8d7ed1a0e77eb4d6374d182ee9d4ddcddbb.
\end{quote}
The two principal endpoints in namespace \texttt{KiralyClutter} are
\begin{quote}\small\ttfamily
counterexample : IsClutter C11 $\land$ $\neg$ IsIdeal C11 $\land$ $\neg$ KiralyProperty C11\\
kiraly\_conjecture\_false : $\neg$ KiralyConjecture.
\end{quote}
The second endpoint universally quantifies over finite ground types and clutters in its negated conjecture. The first supplies the concrete instance. No theorem from a source paper is imported as an extra Lean axiom: the mathematical source attributions concern definitions, question and context; the finite instance and deductions are checked directly.

The portable \texttt{proof-packet.zip} contains the exact theorem, Lean toolchain file, complete nine-package Lake manifest, explanatory JSON, certificate inspector and configurable replay script. It intentionally contains no historical evidence tree, dependency source tree or precompiled theorem object. From a fresh extraction, use
\begin{quote}\small\ttfamily
python3 replay.py --dependency-root /path/to/.lake/packages \\
\hspace*{1em}--lean /path/to/lean-4.34.1/bin/lean --output replay-output
\end{quote}
The dependency root may be selected by environment variable; the script constructs \texttt{LEAN\_PATH} from it and checks the pinned checkout revisions without modifying dependencies. The optional \texttt{--leanchecker} runs a module check after compilation. Paths are passed as argument vectors, not shell interpolation.

Lean is pinned to version 4.34.1. The exact compiler and mathlib revisions are:
\begin{quote}\small\ttfamily
Lean: 5045d0056413266e57c625dcd7c365b10e377c52\\
mathlib: d13f23b723b8a846827a245b89c10fc7d3f11612
\end{quote}
The accompanying replay logs record fresh archive extraction, two-thread compilation with warnings as errors, the printed axiom dependencies of both endpoints, and the optional module check. Only \texttt{propext}, \texttt{Classical.choice} and \texttt{Quot.sound} occur in the endpoint axiom output. There are no admitted proof obligations in these endpoints.

\paragraph{Limits of this replay.}
The run uses the current external prebuilt mathlib/package cache. It is a fresh build of the theorem from its bytes, not a rebuild of Lean and mathlib from network source. Checkout revision pins do not authenticate the user-writable prebuilt dependency objects, which are read-only by convention. The Cli checkout lacks a built library in this cache; that unused path is omitted and reported. The optional \texttt{leanchecker} run checks the new module against imports; it is not a fresh full-mathlib replay. No full dependency rebuild or global literature clearance is claimed. Source PDF hashes, the theorem-to-paper mapping, input verification and artifact manifests are supplied outside the manuscript in the reproducibility dossier.

\section{Consequence and candidate status}
The obstruction proves that a nonideal finite clutter need not admit the proposed choice maps. Since any single-$h$ witness restricts to such maps, it also refutes the implication posed in Pap's report. The consequence uses the full blocker domain in the statement but only certified members in the proof; a successful core-only choice construction cannot establish the disputed implication on its own.

The completed independent fidelity review \cite{Fidelity} reports PASS for its four stated gates; this draft uses that review for source-form mapping and records it as provenance, not as novelty approval. The completed separate significance review \cite{Significance} reports source PASS, mathematical significance PASS and qualified novelty PASS under a bounded-prior standard. Its assessed contribution is the all-assignments obstruction on the full blocker, rather than the known object, nonideality or formalization. It found no demonstrated public subsumer in the investigated sources, but the original L\"utolf--Margot catalogue remains unread and later-literature coverage is incomplete. Neither catalogue chronology nor the historical ``open'' label clears those gaps. No worldwide priority is certified. These reviews did not review the finished manuscript and do not replace final acceptance by the parent investigator. This manuscript and packet grant zero acceptance credit. No publication, submission or contact with third parties is authorized or performed.

\clearpage
\begin{thebibliography}{9}
\raggedright
\bibitem{PapTR} J\'ulia Pap, \emph{A note on a conjecture on clutters}, EGRES Technical Report TR-2010-09, November 2010 version, Definition 1.1 and Conjecture 1.2, p.~1.\newline
\url{https://egres.elte.hu/tr/egres-10-09_November_2010.pdf}.
\bibitem{PapRevised} J\'ulia Pap, \emph{A note on a conjecture on clutters}, EGRES TR-2010-09, November 2010, revised January 2011. Definition 1.1 and Conjecture 1.2, p.~1. Exact revision retrieved by independent reviewer on 1 November 2026 UTC.\newline
\url{https://egres.elte.hu/tr/egres-10-09.pdf}.
\bibitem{PapThesis} J\'ulia Pap, \emph{Integrality, complexity and colourings in polyhedral combinatorics}, Ph.D. thesis, E\"otv\"os Lor\'and University, title page dated April 2012; archived PDF capture 31 December 2015. Conjecture 3.4.5, pp.~70--71. Author-hosted copy, retrieved by independent reviewer on 1 November 2026 UTC: \url{https://papjuli.github.io/assets/theses/julia_pap_phd_thesis.pdf}. Original URL \url{http://www.cs.elte.hu/math/phd_th/Pap_Julia.pdf}; exact archived version:\newline
\url{https://web.archive.org/web/20151231130051id_/http://www.cs.elte.hu/math/phd_th/Pap_Julia.pdf}.
\bibitem{Egres} Egres Open, \emph{Orientation of nonideal clutters}, snapshot of oldid 1782, last modified 2 January 2015, 17:37. Exact cached bytes and hash in the dossier.\newline
\url{http://lemon.cs.elte.hu/egres/internals/index.php5?title=Orientation_of_nonideal_clutters&oldid=1782}.
\bibitem{MPR} Dillon Mayhew, Irene Pivotto and Gordon Royle, \emph{Structure of Cubic Lehman Matrices}, arXiv:1805.07576v2, posted 31 July 2019; PDF title date 1 August 2019. Definition 1.1; Section 5 and Table 1.\newline
\url{https://arxiv.org/abs/1805.07576v2}.
\bibitem{Fidelity} Independent fidelity review, \emph{Candidate 388087}, internal completed review, 31 October 2026 local / 1 November UTC, lane \texttt{kiraly-fidelity-review-r1}, \texttt{REVIEW.md} and \texttt{review.json}. Copies with hashes in the dossier; not a publication or novelty assessment.
\bibitem{Significance} Independent significance and bounded prior-art review, \emph{MathDB 388087}, completed 31 October 2026 local / 1 November UTC, lane \texttt{kiraly-significance-review-r1}, \texttt{REVIEW.md} and \texttt{review.json}. Copies with hashes in the dossier; qualified bounded-prior finding, not acceptance.
\end{thebibliography}
\end{document}
